%%
%% This is file `anleitung.tex'.
%%
%% Copyright (C) 1994-2002 Axel Sommerfeldt
%% 
%% --------------------------------------------------------------------------
%% 
%% It may be distributed and/or modified under the
%% conditions of the LaTeX Project Public License, either version 1.2
%% of this license or (at your option) any later version.
%% The latest version of this license is in
%%    http://www.latex-project.org/lppl.txt
%% and version 1.2 or later is part of all distributions of LaTeX
%% version 1999/12/01 or later.
%% 
\NeedsTeXFormat{LaTeX2e}[1994/12/01]
\documentclass[a4paper]{ltxdoc}
\setlength\parindent{0pt}
\setlength\parskip{\medskipamount}
\usepackage[german]{babel}
%\usepackage{comment}

\errorcontextlines9
\usepackage[bf,debug]{caption2}[2002/08/03]
\setlength\abovecaptionskip{10pt}
\setlength\belowcaptionskip{10pt}
%\setcaptionmargin{1em}

\newenvironment{example}{\list{}{\leftmargin=10pt\rightmargin\leftmargin}\item[]}{\endlist}

% Various forms of argument:
% (taken from ltxguide.cls:)
%\newcommand{\m}[1]{\mbox{$\langle$\it #1\/$\rangle$}}
%\renewcommand{\arg}[1]{{\tt\string{}\m{#1}{\tt\string}}}
%\newcommand{\oarg}[1]{{\tt[}\m{#1}{\tt]}}
% (taken from ltxdoc.cls:)
%\DeclareRobustCommand\cs[1]{\texttt{\char`\\#1}}
%\providecommand\marg[1]{%
%  {\ttfamily\char`\{}\meta{#1}{\ttfamily\char`\}}}
%\providecommand\oarg[1]{%
%  {\ttfamily[}\meta{#1}{\ttfamily]}}
%\providecommand\parg[1]{%
%  {\ttfamily(}\meta{#1}{\ttfamily)}}

\begin{document}

\newcommand{\purerm}[1]{{\upshape\mdseries\rmfamily #1}}
\newcommand{\puresf}[1]{{\upshape\mdseries\sffamily #1}}
\newcommand{\captionpkt}{\puresf{caption2}}
\def\figurename{Beispiel}

\title{Das \captionpkt\ Paket}
\author{Axel Sommerfeldt\\E-Mail:~caption@sommerfeldt.net}
\date{letzte "Anderung: 3.~August 2002}
\maketitle

\begin{abstract}
Das \captionpkt\ Paket bietet einem viele Wege das Erscheinungsbild der Bildunterschriften
und Tabellen"uberschriften den eigenen W"unschen bzw.\ bestimmten Vorgaben anzupassen.
Hierbei wurde viel Wert auf die reibungslose Zusammenarbeit und in die Integration anderer
Dokumentenklassen bzw.\ Pakete gelegt.
\end{abstract}

\section{Hallo Welt!}

Leider ist diese Anleitung (immer) noch alles andere als fertig, dennoch habe ich mich dazu entschlossen, den Kram trotzdem jetzt schon zu ver"offentlichen, denn:

\begin{enumerate}
\item Die in der letzten Version angegebene E-Mail-Adresse funktioniert leider nicht mehr.
\item Die letzte Version hatte keine Lizenz, so da"s hier Kl"arungsbedarf bestand.
(caption2 ist jetzt unter der LPPL, wie die meisten anderen \LaTeXe-Pakete auch.)
\item Die aktuelle Version von subfigure (2.1) machte Anpassungen an \captionpkt\ notwendig.
\item Die letzte Version von \captionpkt\ hatte auch keine brauchbare Anleitung.
\item Und "uberhaupt wurde es endlich mal Zeit f"ur eine Releaseversion.
\end{enumerate}

Zusammenfassend l"a"st sich also sagen, da"s diese Version trotz seiner Defizite immer noch besser ist als die letzte Version auf CTAN.

Ich werde mich bem"uhen, die Anleitung in den n"achsten Wochen fertigzubekommen. Der aktuelle Zwischenstand kann solange "uber%
\begin{quote}\texttt{http://www.sommerfeldt.net/latex/}\end{quote}%
bezogen werden.

Sobald die Anleitung fertig ist, werde ich auch beginnen, Verbesserungen und Erweiterungen in \captionpkt\ einzupflegen.

\section{Einleitung}
Bildunterschriften bzw.\ Tabellen"uberschriften werden von den Standard"=Dokumentenklassen
eher stiefm"utterlich behandelt -- sie werden einfach als ganz normaler Absatz gesetzt.
Dies hat aber diverse Nachteile:

\begin{itemize}
\item Es passiert leicht, da"s man beim Weiterlesen eines
Absatzes auf der n"achsten Seite zun"achst versehentlich in eine Bildunterschrift anstelle
des normalen Textes ger"at. Dies st"ort den Leseflu"s.
\item Sie passen sich u.U.~nicht in das Layout des restlichen Dokumentes ein. Sind z.B. die
"Uberschriften des Dokumentes serifenlos, so w"are es schon aus optischen Gr"unden w"unschenswert,
die Bildunterschriften und Tabellen"uberschriften ebenfalls serifenlos zu erhalten.
\item Es ist nicht ohne Kenntnisse von \LaTeXe -Internata m"oglich, diese Unter- bzw.\ "Uberschriften eigenen W"unschen oder gegebenen Vorlagen anzupassen.
\end{itemize}

All diese Nachteile versucht \captionpkt\ zu beheben. Hierzu werden die unterschiedlichsten
Gestaltungsmerkmale ber"ucksichtigt: R"ander, Abst"ande zum Text, Schriftarten und -gr"o"sen etc.

%\section{Installation}
%Wie auch viele andere \LaTeXe\ Pakete liegt dieses Paket im sog.\ "`docstrip"� Format vor, welches erst
%ausgepackt werden mu"s. Ein \LaTeXe -Lauf der Datei \texttt{caption.ins} liefert die eigentlichen
%Paketdateien \texttt{caption.sty} und \texttt{caption2.sty}. Diese Dateien k"onnen dann einfach in das
%gleiche Verzeichnis wie das eigentliche Dokument gelegt werden --
%praktischer ist jedoch eine Installations ins bestehende
%\TeX -System. Wie dies funktioniert entnehmen sie bitte der Anleitung Ihrer \TeX -Distribution, bei
%Unix-Versionen m"ussen in der Regel die Dateien in das Verzeichnis
%\texttt{\$TEXMF/tex/contrib} kopiert werden,
%ein anschlie"sender Aufruf von \texttt{texhash} vermerkt die Datei anschlie"send im System.
%Verwenden Sie MikTeX unter MS-Windows, so kopieren Sie die Dateien bitte nach
%\texttt{\textbackslash texmf\textbackslash tex\textbackslash latex\textbackslash misc}
%und machen die Datei anschlie"send mit Hilfe des Startmen"ueintrages "`Refresh Filename Database"� dem
%MikTeX-System bekannt.

%\pagebreak
\section{Verwendung des Paketes}
\label{verwendung}

Das Paket wird mittels
%
\begin{example}
|\usepackage|\oarg{options}|{caption2}[2002/08/03]|
\end{example}
%
eingebunden. Bitte verwenden Sie immer die neuere Version |caption2| und nicht mehr die "altere Version |caption|, diese steht nur noch aus Kompatibilit"atsgr"unden zur Verf"ugung.

Folgende Optionen stehen zur Verf"ugung:

%\DeleteShortVerb{\|}%
\small %footnotesize
\begin{tabular}{lll}
%\MakeShortVerb{\|}%
\textbf{Option} & \textbf{entsprechendes Kommando} & \textbf{Abschnitt}\\[0.5ex]
%\hline
|normal|       & |\captionstyle{normal}|           & \ref{captionstyle} \\
|center|       & |\captionstyle{center}|           & \ref{captionstyle} \\
|flushleft|    & |\captionstyle{flushleft}|        & \ref{captionstyle} \\
|flushright|   & |\captionstyle{flushright}|       & \ref{captionstyle} \\
|centerlast|   & |\captionstyle{centerlast}|       & \ref{captionstyle} \\
|hang|         & |\captionstyle{hang}|             & \ref{captionstyle} \\
|indent|       & |\captionstyle{indent}|           & \ref{captionstyle} \\
|scriptsize|   & |\setcaptionfont{\scriptsize}|    & \ref{captionfont} \\
|footnotesize| & |\renewcommand\captionfont{\footnotesize}|  & \ref{captionfont} \\
|small|        & |\renewcommand\captionfont{\small}|         & \ref{captionfont} \\
|normalsize|   & |\renewcommand\captionfont{\normalsize}|    & \ref{captionfont} \\
|large|        & |\renewcommand\captionfont{\large}|         & \ref{captionfont} \\
|Large|        & |\renewcommand\captionfont{\Large}|         & \ref{captionfont} \\
|up|           & |\renewcommand\captionlabelfont{\upshape}|  & \ref{captionfont} \\
|it|           & |\renewcommand\captionlabelfont{\itshape}|  & \ref{captionfont} \\
|sl|           & |\renewcommand\captionlabelfont{\slshape}|  & \ref{captionfont} \\
|sc|           & |\renewcommand\captionlabelfont{\scshape}|  & \ref{captionfont} \\
|md|           & |\renewcommand\captionlabelfont{\mdseries}| & \ref{captionfont} \\
|bf|           & |\renewcommand\captionlabelfont{\bfseries}| & \ref{captionfont} \\
|rm|           & |\renewcommand\captionlabelfont{\rmfamily}| & \ref{captionfont} \\
|sf|           & |\renewcommand\captionlabelfont{\sffamily}| & \ref{captionfont} \\
|tt|           & |\renewcommand\captionlabelfont{\ttfamily}| & \ref{captionfont} \\
|oneline|      & |\onelinecaptionstrue|            & \ref{onelinecaption} \\
|nooneline|    & |\onelinecaptionsfalse|           & \ref{onelinecaption} \\
|float|        & ---                               & \ref{float} \\
|longtable|    & ---                               & \ref{longtable} \\
|ignoreLTcapwidth| & ---                           & \ref{longtable} \\
|subfigure|    & ---                               & \ref{subfigure} \\
|all|          & ---                               & \ref{packages} \\
|none|         & ---                               & \ref{packages} \\
\end{tabular}
\normalsize

\pagebreak
\section{Vordefinierte Stile}
\label{captionstyle}
%
\begin{macro}{\captionstyle}
Neben der M"oglichkeit neue Stile zu definieren bietet \captionpkt\ die vordefinierten Typen |normal|, |center|, |flushleft|, |flushright|, |centerlast|, |hang| sowie |indent|. Zwischen diesen kann man im Dokument mittels des Kommandos |\captionsstyle| umschalten, z.B.~wird mittels des Kommandos
%
\begin{example}
|\captionstyle{center}|
\end{example}
%
auf den Stil |center| umgeschaltet. Geschieht dies au"serhalb einer Abbildung, so geschieht dies global f"ur alle nachfolgenden Abbildungen, geschieht dies hingegen innerhalb einer Abbildung, so geschieht dies nur lokal f"ur die aktuelle Abbildung.
\end{macro}

\subsection*{Der Stil |normal|}

Dies ist der Stil, den wir von den Standard \LaTeXe-Klassen |article|, |report| und |book| gewohnt sind. \captionpkt\ setzt diesen Typ daher als Standardtyp.

\captionstyle{normal}
\captionof{figure}[...]{Dies ist eine Unterschrift des Stils |normal|. Dies ist eine Unterschrift des Stils |normal|. Dies ist eine Unterschrift des Stils |normal|.}

\subsection*{Der Stil |center|}

Hier wird die Unterschrift innerhalb der \LaTeXe-Umgebung |center| gesetzt, jede Zeile wird also zentriert gesetzt.

\captionstyle{center}
\captionof{figure}[...]{Dies ist eine Unterschrift des Stils |center|. Dies ist eine Unterschrift des Stils |center|. Dies ist eine Unterschrift des Stils |center|.}

\subsection*{Der Stil |flushleft|}

Hier wird die Unterschrift innerhalb der \LaTeXe-Umgebung |flushleft| gesetzt, jede Zeile wird also linksb"undig gesetzt.

\captionstyle{flushleft}
\captionof{figure}[...]{Dies ist eine Unterschrift des Stils |flushleft|. Dies ist eine Unterschrift des Stils |flushleft|. Dies ist eine Unterschrift des Stils |flushleft|.}

\subsection*{Der Stil |flushright|}

Hier wird die Unterschrift innerhalb der \LaTeXe-Umgebung |flushright| gesetzt, jede Zeile wird also rechtsb"undig gesetzt.

\captionstyle{flushright}
\captionof{figure}[...]{Dies ist eine Unterschrift des Stils |flushright|. Dies ist eine Unterschrift des Stils |flushright|. Dies ist eine Unterschrift des Stils |flushright|.}

\subsection*{Der Stil |centerlast|}

Hier wird die Unterschrift bis auf die letzte Zeile normal, also als Blocksatz gesetzt, die letzte Zeile erscheint hingegen zentriert.

\captionstyle{centerlast}
\captionof{figure}[...]{Dies ist eine Unterschrift des Stils |centerlast|. Dies ist eine Unterschrift des Stils |centerlast|. Dies ist eine Unterschrift des Stils |centerlast|.}

\subsection*{Der Stil |hang|}

Hier wird die Unterschrift so gesetzt, da"s jede Zeile unterhalb der anderen "`h"angt"'.

\captionstyle{hang}
\captionof{figure}[...]{Dies ist eine Unterschrift des Stils |hang|. Dies ist eine Unterschrift des Stils |hang|. Dies ist eine Unterschrift des Stils |hang|.}

\subsection*{Der Stil |indent|}
%
%\begin{decl}
%|\captionindent|
%\end{decl}
%
Dieser Typ "ahnelt dem Typ |hang|, allerdings ist hier der Einzug der unteren Zeilen nicht durch die L"ange der Bezeichnung bestimmt, sondern wird durch die L"ange |\captionindent| festgelegt. W"unscht man also einen Einzug von 1\,cm, so wird dies mittels
%
\begin{example}
|\setlength{\captionindent}{1cm}|
\end{example}
%
bewerkstelligt. Der Einzug ist auf null voreingestellt, ohne explizite Angabe unterscheidet sich dieser Stil also nicht vom Stil |normal|.

\captionstyle{indent}
\setlength\captionindent{1cm}
\captionof{figure}[...]{Dies ist eine Unterschrift des Stils |indent|. Dies ist eine Unterschrift des Stils |indent|. Dies ist eine Unterschrift des Stils |indent|.}

\subsection{Anpassungsm"oglichkeiten der vordefinierten Typen}
%\label{predefined}
\captionstyle{normal}

Die vordefinierten Typen lassen sich anpassen, viele Parameter lassen sich ver"andern, wie etwa der Rand, die Breite, die Schrifttypen.

\subsubsection{Anpassung der Schrifttypen}
\label{captionfont}

Die Unterschriften werden bei den vordefinierten Typen immer wie folgt gesetzt:
%
\begin{example}
|{\captionfont|\\
\hspace*{1em}|{\captionlabelfont |\meta{Bezeichnung}|\captionlabeldelim}%|\\\hspace{1em}
\hspace*{1em}|\captionlabelsep |\meta{Unterschrift}|}|
\end{example}
%
\cs{captionfont}, \cs{captionlabelfont}, \cs{captionlabeldelim} und \cs{captionlabelsep} sind dabei wie folgt vorbelegt:

\begin{tabular}{ll}
\textbf{Parameter}   & \textbf{Vorbelegung}\\[0.5ex]
|\captionfont|       & |{}|\\
|\captionlabelfont|  & |{}|\\
|\captionlabeldelim| & |{:}|\\
|\captionlabelsep|   & |{|\textvisiblespace|}|
\end{tabular}

M"ochte man also z.B.~die Unterschriften generell etwas kleiner als den normalen Text und ferner die Bezeichnung in Italics gesetzt, so l"a"st sich dies durch
%
\begin{example}
|\usepackage[small,it]{caption2}|
\end{example}
%
in der Preambel des Dokuments oder innerhalb des Dokuments durch
%
\begin{example}
|\renewcommand\captionfont{\small}|\\
|\renewcommand\captionlabelfont{\itshape}|
\end{example}
%
bewerkstelligen.
%
\normalcaptionparams
\renewcommand\captionfont{\small}
\renewcommand\captionlabelfont{\itshape}
\captionof{figure}[...]{Dies ist eine Unterschrift}
%
M"ochte man hingegen den Text selber, aber nicht die Bezeichnung in Italics gesetzt haben, so geht dies (etwas umst"andlicher) wie folgt:
%
\begin{example}
|\renewcommand*\captionfont{\small\itshape}|\\
|\renewcommand*\captionlabelfont{\upshape}|
\end{example}
%
\normalcaptionparams
\renewcommand\captionfont{\small\itshape}
\renewcommand\captionlabelfont{\upshape}
\captionof{figure}[...]{Dies ist eine Unterschrift}

\subsubsection{Anpassung der Breite bzw.\ des Randes}
\label{captionwidth}

Die Breite der gesammten Unterschrift kann explizit per |\setcaptionwidth| gesetzt werden, um etwa zu erreichen, da"s die Breite der Unterschrift mit der Breite der Abbildung identisch ist. Die Unterschrift selber wird hierbei zentriert.

\begin{example}
|\setcaptionwidth{.5\textwidth}|\\
|\caption{Diese Unterschrift ist lediglich|\\
|         die halbe Textbreite breit.}|
\end{example}

\normalcaptionparams
\setcaptionwidth{.5\textwidth}
\captionof{figure}{Diese Unterschrift ist lediglich die halbe Textbreite breit.}

Alternativ kann auch der Rand explizit mit |\setcaptionmargin| gesetzt werden, das obenstehende Beispiel lie"se sich also auch so realisieren:

\begin{example}
|\setcaptionmargin{.25\textwidth}|\\
|\caption{Diese Unterschrift ist lediglich|\\
|         die halbe Textbreite breit.}|
\end{example}

\normalcaptionparams
\setcaptionmargin{.25\textwidth}
\captionof{figure}{Diese Unterschrift ist lediglich die halbe Textbreite breit.}

\subsubsection{Sonstige Anpassungen}
\label{onelinecaption}
%
\LaTeX\ mi"st vor dem Setzen der Unterschrift die Breite dieser -- ist sie kleiner oder gleich der Textbreite, so gehen alle manuell gesetzten Umbr"uche etc. verloren und die Unterschrift wird einzeilig zentriert gesetzt. Ist sie jedoch l"anger, wird sie normal als Absatz gesetzt. Das \captionpkt\ bildet diesen Mechanismus nach, er kann jedoch mit |\onelinecaptionsfalse| (oder der Paketoption |nooneline|) ab- bzw.\ mit |\onelinecaptionstrue| (oder der Paketoption |oneline|) angeschaltet werden.
%
\begin{example}
|\onelinecaptionstrue|\\
|\caption{Zeile 1\\Zeile 2}|
\end{example}
%
\normalcaptionparams
\onelinecaptionstrue
\caption{Zeile 1\\Zeile 2}
%
\begin{example}
|\onelinecaptionsfalse|\\
|\caption{Zeile 1\\Zeile 2}|
\end{example}
%
\normalcaptionparams
\onelinecaptionsfalse
\captionof{figure}{Zeile 1\\Zeile 2}
%
\subsubsection{Sonstige Features}
Ansonsten kann \cs{caption*} auch nachgebildet werden:
\begin{example}
|\captionlabelfalse|\\
|\caption|\oarg{xxx}\marg{xxx}
\end{example}
%
\subsubsection{Obsolete Parameter}
Der Parameter |\captionsize| ist obsolete. Er wird zwar noch -- soweit m"oglich -- emuliert, sollte aber nicht mehr verwendet werden.

\section{Selbstdefinierte Stile}
Wem die M"oglichkeiten bis jetzt nicht ausreichten, der kann auch selber eigene Unterschriftentypen definieren. Innerhalb dieser kann man bei Bedarf auch die M"oglichkeiten, die die vordefinierten Stile bieten, einbetten.

\subsection{Deklaration von eigenen Stilen}

Analog zu dem \TeX-Kommando |\def| bzw.\ den \LaTeX-Kommandos |\newcommand| und |\renewcommand| werden die Kommandos |\defcaptionstyle|, |\newcaptionstyle| und |\renewcaptionstyle| zur Verf"ugung gestellt. Die Syntax ist bei allen 3 Kommandos die gleiche:

\begin{quote}
|\defcaptionstyle|\marg{Name}\marg{Kommandos}\\
|\newcaptionstyle|\marg{Name}\marg{Kommandos}\\
|\renewcaptionstyle|\marg{Name}\marg{Kommandos}
\end{quote}

Innerhalb des eigenen Stils stehen die Bezeichnung als |\captionlabel| sowie der eigentliche Unterschriftentext als |\captiontext| zur Verf"ugung. Wollte man also z.B. die Unterschrift als ganz normalen Absatz setzen, so lie"se sich dies so realisieren:
%
\begin{example}
|\newcaptionstyle{absatz}{\captionlabel: \captiontext\par}|\\
...\\
|\captionstyle{absatz}|\\
|\caption{Dies ist eine Unterschrift.}|
\end{example}
%
\newcaptionstyle{absatz}{\captionlabel: \captiontext\par}
\captionstyle{absatz}
\captionof{figure}{Dies ist eine Unterschrift.}
%
Dem Spieltrieb sind hierbei keine Grenzen gesetzt:
%
\begin{example}
|\newcaptionstyle{fancy}{\textsf{\captionlabel}\\\captiontext\par}|\\
...\\
|\captionstyle{fancy}|\\
|\caption{Dies ist eine Unterschrift.}|
\end{example}
%
\newcaptionstyle{fancy}{\textsf{\captionlabel}\\\captiontext\par}
\captionstyle{fancy}
\captionof{figure}{Dies ist eine Unterschrift.}
%
...oder...
%
\begin{example}
|\newcaptionstyle{fancy2}{\captiontext\hfill\textit{(\captionlabel)}}|\\
...\\
|\captionstyle{fancy2}|\\
|\caption{Dies ist eine Unterschrift.}|
\end{example}
%
\newcaptionstyle{fancy2}{\captiontext\hfill\textit{(\captionlabel)}}
\captionstyle{fancy2}
\captionof{figure}{Dies ist eine Unterschrift.}

\subsection{Einbindung vorhandener Hilfsmittel}

\begin{quote}
|\usecaptionstyle|\marg{Name des Stils}\\
|\normalcaptionparams|
\end{quote}
%
Bei der Definition von eigenen Stilen kann ggf.\ aus dem Fundus von \captionpkt\ zur"uckgegriffen werden. M"ochte man z.B.~auf einem vorhandenen Stil aufsetzen, diesen also lediglich modifiziert anbieten, so kann dies innerhalb der Definition mittels |\usecaptionstyle| geschehen. Weiterhin wird das Kommando |\normalcaptionparams| bereitgestellt, mit dem alle Einstellungen der vordefinierten Stile lokal(!) zur"uckgesetzt werden, so da"s man auf einen wohldefinierten Zustand aufsetzen kann. M"ochte man z.B.~einen Stil definieren, der dem Stil |centerlast| entspricht, allerdings immer als Absatz ohne R"ander gesetzt wird und mit einem Punkt anstelle des Doppelpunktes und einer fetten Beschriftung, so l"asst sich dies so realisieren:
%
\begin{example}
|\newcaptionstyle{mystyle}{%|\\
|  \normalcaptionparams|\\
|  \renewcommand\captionlabelfont{\bfseries}%|\\
|  \renewcommand\captionlabeldelim{.}%|\\
|  \onelinecaptionsfalse|\\
|  \usecaptionstyle{centerlast}}|\\
...\\
|\captionstyle{mystyle}|\\
|\caption{Dies ist eine sehr sehr sehr sehr sehr sehr|\\
|         sehr sehr sehr sehr sehr sehr lange Unterschrift.}|
\end{example}
%
\newcaptionstyle{mystyle}{%
  \normalcaptionparams
  \renewcommand\captionlabelfont{\bfseries}%
  \renewcommand\captionlabeldelim{.}%
  \onelinecaptionsfalse
  \usecaptionstyle{centerlast}}
\captionstyle{mystyle}
\captionof{figure}{Dies ist eine sehr sehr sehr sehr sehr sehr
         sehr sehr sehr sehr sehr sehr lange Unterschrift.}
%
Ein weiteres Beispiel, welches einen neuen Stil |hangandleft| definiert. Dieser unterscheidet sich vom Stil |hang| derart, da"s der eigentliche Text linksb"undig gesetzt wird:
%
\begin{example}
|\newcaptionstyle{hangandleft}{%|\\
|  \let\oldcaptiontext\captiontext|\\
|  \def\captiontext{\raggedright\oldcaptiontext}%|\\
|  \usecaptionstyle{hang}}|\\
|\captionstyle{hangandleft}|\\
|\caption{Dies ist eine sehr sehr sehr sehr sehr sehr|\\
|         sehr sehr sehr sehr sehr sehr sehr sehr sehr|\\
|         sehr sehr sehr sehr sehr sehr sehr sehr sehr|\\
|         sehr sehr sehr sehr sehr sehr lange Unterschrift.}|
\end{example}

\newcaptionstyle{hangandleft}{%
  \let\oldcaptiontext\captiontext
  \def\captiontext{\raggedright\oldcaptiontext}%
  \usecaptionstyle{hang}}
\captionstyle{hangandleft}
\captionof{figure}{Dies ist eine sehr sehr sehr sehr sehr sehr
         sehr sehr sehr sehr sehr sehr sehr sehr sehr
         sehr sehr sehr sehr sehr sehr sehr sehr sehr
         sehr sehr sehr sehr sehr sehr lange Unterschrift.}
%
Aber auch f"ur Stile, die nicht auf anderen Stilen aufsetzen, gibt es Hilfsmittel:
%
\begin{quote}
|\onelinecaption|\marg{Einzeilige Unterschrift}\marg{Mehrzeilige Unterschrift}
\end{quote}
%
|\onelinecaption| bildet den \LaTeX-Mechanismus nach, einzeilige Unterschriften zentriert zu setzen. Ist die Unterschrift \meta{Einzeilige Unterschrift}, also der erste Parameter, kleiner oder gleich der Textbreite, so wird diese zentriert gesetzt. Anderenfalls wird die \meta{Mehrzeilige Unterschrift} verwendet. Die Standarddefinition von \LaTeX\ f"ur eine Unterschrift lie"se sich also wie folgt nachbilden:
%
\begin{example}
|\newcaptionstyle{latex}{%|\\
|  \onelinecaption|\\
|    {\captionlabel: \captiontext}%|\\
|    {\captionlabel: \captiontext\par}}|\\
...\\
|\captionstyle{latex}|\\
|\caption{Dies ist eine kurze Unterschrift.}|\\
\%| bzw.|\\
|\caption{Dies ist eine sehr sehr sehr sehr sehr sehr|\\
|         sehr sehr sehr sehr sehr sehr lange Unterschrift.}|
\end{example}
%
\newcaptionstyle{latex}{%
  \onelinecaption
    {\captionlabel: \captiontext}%
    {\captionlabel: \captiontext\par}}
\captionstyle{latex}
\onelinecaptionstrue
\captionof{figure}{Dies ist eine kurze Unterschrift.}
\captionof{figure}{Dies ist eine sehr sehr sehr sehr sehr sehr
         sehr sehr sehr sehr sehr sehr lange Unterschrift.}
%
|\onelinecaption| beachtet hierbei das boolische Flag |\onelinecaptions|, m"ochte man dies nicht, so mu"s man dies (logischerweise) durch ein |\onelinecaptionstrue| unmittelbar vor der Verwendung von |\onelinecaption| ausschalten, also etwa so:
%
\begin{example}
|\newcaptionstyle{strictlatex}{%|\\
|  \onelinecaptionstrue|\\
|  \onelinecaption|\\
|    {\captionlabel: \captiontext}%|\\
|    {\captionlabel: \captiontext\par}}|
\end{example}
%
%\begin{decl}
%|\usecaptionmargin|
%\end{decl}
%
|\usecaptionmargin| setzt den zuvor mittels |\setcaptionmargin| oder |\setcaptionwidth| gesetzten Rand. Hierzu ein komplexeres Beispiel, welches neben dem gesetzten Rand auch |\onelinecaption| verwendet sowie die gesetzten Schrifttypen beachtet:
%
\begin{example}
|\newcommand*\centerlast{%|\\
|  \addtolength{\leftskip}{0pt plus 1fil}%|\\
|  \addtolength{\rightskip}{0pt plus -1fil}%|\\
|  \setlength{\parfillskip}{0pt plus 2fil}}|\\
|\newcaptionstyle{fancy}{%|\\
|  \usecaptionmargin\captionfont|\\
|  \onelinecaption|\\
|    {{\captionlabelfont\captionlabel\captionlabeldelim}%|\\
|     \captionlabelsep\captiontext}%|\\
|    {{\centering\captionlabelfont\captionlabel\par}%|\\
|      \centerlast\captiontext\par}}|\\
...\\
|\captionstyle{fancy}|\\
|\caption{Dies ist eine kurze Unterschrift.}|\\
|\caption{Dies ist eine sehr sehr sehr sehr sehr sehr|\\
|         sehr sehr sehr sehr sehr sehr lange Unterschrift.}|
\end{example}
%
\newcommand*\centerlast{%
  \addtolength{\leftskip}{0pt plus 1fil}%
  \addtolength{\rightskip}{0pt plus -1fil}%
  \setlength{\parfillskip}{0pt plus 2fil}}
\renewcaptionstyle{fancy}{%
  \usecaptionmargin\captionfont
  \onelinecaption
    {{\captionlabelfont\captionlabel\captionlabeldelim}%
     \captionlabelsep\captiontext}%
    {{\centering\captionlabelfont\captionlabel\par}%
      \centerlast\captiontext\par}}
\captionstyle{fancy}
\normalcaptionparams
\renewcommand\captionlabelfont{\bfseries}
\captionof{figure}{Dies ist eine kurze Unterschrift.}
\captionof{figure}{Dies ist eine sehr sehr sehr sehr sehr sehr
         sehr sehr sehr sehr sehr sehr lange Unterschrift.}

%\dummycaptionstyle

\section{Sonstiges}
\captionstyle{normal}
\normalcaptionparams

\begin{quote}
|\captionof|\marg{Floattyp}\oarg{Kurzform}\marg{Langform}
\end{quote}
%
Mittels des Kommandos |\captionof| lassen sich Unterschriften innerhalb des normalen Textes, also au"serhalb von gleitenden Umgebungen, setzen. Ein Beispiel:
%
\begin{example}
|\captionof{table}{Dies ist keine Tabelle}|\\
\end{example}
%
\captionof{table}{Dies ist keine Tabelle}

\begin{quote}
|\setlength{\abovecaptionskip}|\marg{Wert}\\
|\setlength{\belowcaptionskip}|\marg{Wert}
\end{quote}
%
Die L"angen |\abovecaptionskip| und |\belowcaptionskip| beeinhalten den Abstand "uber bzw.\ unter der eigentlichen Unterschrift. In den \LaTeX-Dokumentenklassen \textsf{article}, \textsf{report} und \textsf{book} wird |\abovecaptionskip| auf 10pt voreingestellt, |\belowcaptionskip| auf 0pt.

%\section{Interaction mit andere Klassen}
%\subsection{ucthesis}

\section{Interaktion mit anderen Paketen}

\captionpkt\ erkennt ggf. geladene Pakete und modifiziert diese so, da"s sich die Stil"anderungen auch automatisch auf diese Pakete auswirken. Dies wird meist durch sog.\ Pseudo-Stile realisiert, diese Stile werden zwar definiert, verwenden jedoch -- ggf.\ nach der Modifikation einiger Parameter -- den aktuell eingestellten Stil.

Diese Pseudo-Stile k"onnen daher nicht mittels |\captionstyle| ausgew"ahlt werden, denn ansonsten w"aren sie selbst ja der aktuell eingestellte Stil. Diese Pseudo-Stile k"onnen bei Bedarf umdefiniert werden, m"ochte man z.B.\ unabh"angig von der akt.\ Stilauswahl in Verbindung mit dem Paket \puresf{longtable} immer den Stil |normal| angewendet wissen, so lie"se sich das mit
%
\begin{example}
|\renewcaptionstyle{longtable}{\usecaptionstyle{normal}}|
\end{example}
%
\subsection{Das \puresf{float} Paket}
\label{float}
%\subsection{koma-script}

\subsection{Das \puresf{longtable} Paket}
\label{longtable}

Wird \puresf{longtable} $3.15$ oder neuerer verwendet, wird automatisch ein Pseudo-Stil |longtable| angelegt, welcher innerhalb der |longtable|-Umgebung verwendet wird.

%\subsection{ntgclass}

%\subsection{rotating}
%\label{rotating}

\subsection{subfigure}
\label{subfigure}

Inkompatibel bei `kurzen' Bildunterschriften:

1. |hang,centerlast,nooneline|
2. |hang,center,nooneline|

%\subsection{ucthesis}

%\section{Von caption nach caption2}

% This new version is nearly compatible with the lastest official release
% (version 1.4b), so you can use the old documentation so far.
% Here is what differs this version from version 1.4b:
%
% \begin{itemize}
% \item
%   If the caption package will detect a loaded float package, it will
%   \emph{not} redefine the boxed style of floats anymore. If you want to
%   have the old behaviour, you have to specify the new option |boxed| to
%   the caption2 package.
% \item
%   Anything said about the subfigure package in the old doc isn't
%   true anymore; the caption package is now adapted to the new version
%   2.0 of this package. Especially the caption package will \emph{not}
%   redefine |\@thesubfigure| and |\@thesubtable| anymore and it will
%   \emph{not} set |\subcapsize| -- you have to do this now for yourself
%   if you want to, e.g.\ with the following code:
%   \begin{example}
%       |\usepackage[normalsize]{subfigure}|\\
%       |\usepackage[large]{caption}|
%   \end{example}
%
%   So you can load the caption2 package \emph{before} loading the subfigure
%   package now without problems, in fact this is recommend now. Don't care
%   about what the old doc or the doc of the subfigure package is telling you!
% \end{itemize}
%
% As a summary, the new caption package won't lead into different results
% of your documents just because of loading it (without options).
%
% If you use the new command |\setcaptionwidth| to set the absolut width of a
% caption, you are not allowed to change |\captionmargin| anymore!
% Instead, use the new command |\setcaptionmargin| to do this.
%
% Longtables will still take care of |\LTcapwidth|, even if you are setting
% your own width via |\setcaptionwidth| or |\setcaptionmargin|. To get rid
% of this, use the following code just after loading the caption2 package:
% \begin{example}
%   |\dummycaptionstyle{longtable}{}|
% \end{example}
% or just specify the new package option |longtable|.

\section{Bekannte Probleme}

ucthesis:
1. Korrektur \cs{abovecaptionskip} bzw. \cs{belowcaptionskip} klappt hier nicht
(Per Hand: Bei topcaptions \cs{abovecaptionskip} auf 0pt setzen)
2. In die Definition von \cs{captionfont} mu"s ein \cs{ssp} eingebaut werden,
   also im einfachsten Falle:
  |\renewcommand\captionfont{\ssp}|

\section{Andere n"utzliche Pakete}

float
rotating
sidecap
subfigure

\section{Literaturhinweise}

Folgende Dokumente m"oechte ich an dieser Stelle jedem ans Herz legen:
%
\begin{itemize}
\item Die DANTE-FAQ:\begin{quote}\texttt{http://www.dante.de/faq/de-tex-faq/}\end{quote}
\item\textsf{epslatex} von Keith Reckdahl enth"alt viele n"utzliche Tips im Zusammenhang mit der Einbindung von Graphiken in \LaTeXe. Das Dokument ist unter \begin{quote}\texttt{ftp://ftp.dante.de/pub/tex/info/}\end{quote} als \texttt{epslatex.ps} bzw. \texttt{epslatex.pdf} zu finden.
\end{itemize}

\section{FAQ - oft gestellte Fragen}

\subsubsection*{Ich m"oechte eine leere Unterschrift erzeugen, wie bekomme ich den st"orenden Doppelpunkt weg?}
Indem man sowohl \cs{captionlabeldelim} als auch \cs{captionlabelsep} auf nichts definiert:
%
\begin{example}
|\renewcommand\captionlabeldelim{}|\\
|\renewcommand\captionlabelsep{}|\\
|\caption{}|
\end{example}
%
wird zu:
%
\renewcommand\captionlabeldelim{}
\renewcommand\captionlabelsep{}
\captionof{figure}{}
%
\section{Was tun bei Problemen?}
Steht man vor einem Problem bei der Verwendung des \captionpkt\ Paketes, so bitte ich um eine kurze Nachricht mit einer pr"azisen Angabe des Problems. Dieser Nachricht sollte unbedingt ein m"oglichst kurzes, "ubersetzbares(!) \LaTeXe-Dokument sowie die dazugeh"'orige \texttt{log}-Datei beiliegen, wo der Fehler bzw.\ das Problem auftritt, damit ich es nachstellen und beheben kann.

Auch Anregungen bzw.\ Fragen zu dieser Anleitung sind herzlich willkommen, ich kann nur dann etwas verbessern, wenn ich wei"s, da"s es etwas verbesserungsw"urdiges gibt. Danke!

%\section{Danksagungen}

% I would like to thank David Carlisle for his help writing the longtable
% support; without the changes in his package this wouldn't become possible.

\end{document}
